\subsection{The sharp dimension as a function of field degree}
Order the multiset containing $n$ copies of $1/p$ for each rational prime $p$ in decreasing order and let $S_n(d)$ be the sum of its first $d$ terms.
Among all number fields of degree $n$, the least number of prime-ideal coordinates in which a two-deletion edit strictly improves the greedy cube is
\[
 d_n=\min\{d:S_n(d)>11/2\}.
\]
Every field of degree $n$ requires at least this many coordinates for such an edit. At least one field of degree $n$ attains this dimension. The minimum is finite by divergence of the prime reciprocal sum. The invariant measures exactly the class of edits classified above.

Indeed, for marker weight $w$ the edge formula requires $w^2(1+\sum_i u_i)>1+w$ for a strict gain, since $u_i\le 1/2$ rules out gain from a single spoke. Each satellite contributes $u_i\le w_i$, with equality at satellite bit one. Thus the sum of all the prime-ideal weights must exceed $w+(1+w)/w^2-1$, a decreasing function on $(0,1/2]$ with minimum $11/2$. There are at most $n$ primes above each rational $p$, each of norm at least $p$. Consequently $S_n(d)$ dominates the sum of any $d$ weights in the field.

For the converse, choose a degree-$n$ field in which every rational prime occurring in the first $d_n$ multiset terms splits completely. Such a field exists by coefficientwise Chinese remaindering. Set $g(X)=\prod_{j=1}^n(X-j)$ and, at every specified prime $p$, demand that a monic degree-$n$ polynomial $f$ be congruent to $g$ modulo $p^{2M+1}$ with $M=\max_j v_p(g'(j))$. Hensel's lemma at each $j$ gives a root whose difference from $j$ has valuation $>M$, and the resulting roots are distinct since $v_p(j-k)\le M$ for $j\ne k$. Impose Eisenstein's criterion at one further prime to guarantee irreducibility. Then $\mathbb Q[X]/(f)$ has the required degree and splitting by its tensor products with $\mathbb Q_p$. In these coordinates choose a norm-two marker and take all satellite bits equal to one. The sum condition is sufficient.

\begin{center}
\begin{tabular}{c|rrrrrrrrrr}
$n$&2&3&4&5&6&7&8&9&10&$\ge11$\\ \hline
$d_n$&30\,948&85&27&18&15&14&13&13&12&12
\end{tabular}
\end{center}
For $n\ge12$ the first twelve terms are all $1/2$. For $n=2$ the last prime is $169751$. Integer interval certificates and exact short sums are reproduced by \texttt{check\_field\_thresholds.py}. The case $n=1$ is covered by the formula without requiring enumeration of its prime prefix.

Over the rational integers the threshold nevertheless has a narrow certified scale,
\[
                   3.46\cdot10^{79}<d_1<3.52\cdot10^{79}.
\]
Indeed put $x_-=6.51\cdot10^{81}$ and $x_+=6.57\cdot10^{81}$. The absolute error in replacing the prime reciprocal sum through $x$ by $\log\log x+B$ is at most $1/(2\log^2 x)$ at both endpoints, by Theorem~5 of J.~B.~Rosser and L.~Schoenfeld, \href{https://doi.org/10.1215/ijm/1255631807}{\emph{Approximate formulas for some functions of prime numbers}} (1962). Here $0.261496<B<0.261499$ for the Meissel--Mertens constant. Substitution places the two prime sums on opposite sides of $11/2$. Theorem~1 of the same paper bounds $\pi(x)$ below and above by $(x/\log x)(1+c/\log x)$ with $c=1/2$ and $3/2$, respectively, in this range. Applying these bounds to $\pi(x_-)<d_1\le\pi(x_+)$ gives the interval. The accompanying checker also bounds $B$ directly by truncating $B=\gamma-\sum_p\sum_{k\ge2}1/(kp^k)$ with explicit tail errors.
